\documentclass[11pt]{article}
\usepackage[a4paper,margin=26mm]{geometry}
\usepackage{amsmath,amssymb,amsthm,mathtools,hyperref,booktabs,xurl}
\setlength{\emergencystretch}{2em}
\newtheorem{theorem}{Theorem}
\newtheorem{lemma}[theorem]{Lemma}
\newtheorem{corollary}[theorem]{Corollary}
\newcommand{\Q}{\mathbb Q}\newcommand{\Z}{\mathbb Z}
\newcommand{\A}{\mathcal A}
\newcommand{\Mat}{\operatorname{Mat}}
\newcommand{\Spec}{\operatorname{Spec}}
\title{Shared integer certificates and rational undecidability\\in 99, 117, and 69 variables}
\author{\(\mathbb Z\mathbb Q\) Diophantine Project}
\date{23 November 2026}
\begin{document}\maketitle
\begin{abstract}
We construct simultaneous integer certificates by sharing seven elliptic
point roles, one contact parameter, and one finite-algebra encoder. The
resulting input projections are $\mathbb Z^k$ in $9k+16$ coordinates for
the general source and $11k+16$ for its fully specified version. A separate
ordinary-$abc$ construction uses $7k+4$ coordinates. The mixed-integer
universality theorem then gives rational undecidability bounds $99$, $117$,
and $69$, respectively. The first two use no additional $abc$ hypothesis;
all retained non-pole arithmetic premises of the project sources remain
explicit. The existing Bhargava--Burungale product-pole theorem supplies
the pole converse. These are unrestricted-degree results.

We give the complete ordered simultaneous-witness proofs, the zero-input
and zero-terminal-root cases, and an explicit guarded finite-algebra
encoder valid even for nonreduced algebras. Complete coefficient circuits,
coordinate maps and exact degree ledgers accompany the proof. A separate
quartic source audit is not used in these reductions.
\end{abstract}

\section{Statement}
Put $R=\Z[T_1,\ldots,T_b]$. Let $\A$ be a commutative unital $R$-algebra
free of rank $n\ge1$ as an $R$-module. Fix a basis, elements
$\theta_1,\ldots,\theta_m\in\A$ generating
$\A\otimes_{\mathbb Z}\mathbb Q$ as a $\mathbb Q[T_1,\ldots,T_b]$-algebra, and a guard
$g\in\A$. For a rational specialization $t\in\Q^b$, write
$\A_t=\A\otimes_{R,t}\Q$. A character means a unital algebra homomorphism.

\begin{theorem}\label{main}
An explicit polynomial $F\in R[U]$ can be constructed such that
\begin{equation}\label{equivalence}
 \exists u\in\Q\quad F(t,u)=0
 \quad\Longleftrightarrow\quad
 \exists\phi:\A_t\longrightarrow\Q\quad\phi(g)\ne0
\end{equation}
for every $t\in\Q^b$. Neither reducedness nor etaleness of $\A_t$ is
assumed. In particular, zero roots in any number of quadratic stages are
permitted. If $d$ bounds the total degree of the multiplication-matrix
entries of $g,g\theta_1,\ldots,g\theta_m$, set
\[
 \ell=m+1,\qquad J=\ell(n+1)^2.
\]
One may take
\begin{equation}\label{degree}
 \deg F\le
 2n(J+1)\bigl[1+d\{1+2J(1+\ell(J+1))\}\bigr].
\end{equation}
The bound measures degree in all base coordinates and $U$.
\end{theorem}

The guard may be a product. It then requires each of its factors to be
nonzero at the selected character, without imposing nonvanishing at
unselected characters. No variables beyond $U$ are used in this theorem.

\section{Spectra, including nonreduced algebras}

\begin{lemma}\label{spectrum}
Let $A$ be a finite-dimensional commutative $\Q$-algebra. The eigenvalues
of multiplication by $a\in A$, over $\overline{\Q}$, are the values
$\sigma(a)$ of its geometric characters, with positive multiplicities.
Consequently, for two copies of its multiplication matrix, the eigenvalues
of $M_a\otimes I-I\otimes M_a$ are all differences
$\sigma(a)-\tau(a)$, with multiplicities.
\end{lemma}
\begin{proof}
The commuting multiplication operators can be simultaneously triangularized
over an algebraically closed field. Each diagonal entry is additive,
multiplicative and sends the identity to one, and so is a character.
Alternatively, decompose $A\otimes\overline{\Q}$ as a product of local
Artinian algebras; on each factor multiplication by $a$ is its residue
value plus a nilpotent operator. The tensor difference of upper triangular
matrices has the stated diagonal entries. Neither argument requires that
the nilpotent operators vanish.
\end{proof}

\section{A completely specified polynomial}
Define
\[
 \gamma_0=g,\qquad \gamma_i=g\theta_i\quad(1\le i\le m),
 \qquad A_i=M_{\gamma_i}\in\Mat_n(R).
\]
Set $N=n+1$ and adjoin an extra zero eigenvalue:
\[
 \widetilde A_i=\begin{pmatrix}A_i&0\\0&0\end{pmatrix},\qquad
 D_i=\widetilde A_i\otimes I_N-I_N\otimes\widetilde A_i.
\]
The extra eigenvalue is used to separate every permitted character from
zero. Compute the single monic polynomial
\begin{equation}\label{differencepolynomial}
 \chi(Z)=\prod_{i=0}^{m}\det(ZI_{N^2}-D_i)
       =Z^J+a_1Z^{J-1}+\cdots+a_J,
 \qquad a_0=1.
\end{equation}
Its roots comprise every difference of two character values of a
$\gamma_i$, and every such value itself and its negative.

The following are polynomials in the base coordinates:
\begin{align}
 R_0&=2+N+\sum_{i=0}^{m}\sum_{v,w=1}^{n}(A_i)_{vw}^2
                         +\sum_{j=1}^{J}a_j^2,\label{Rbound}\\
 B_0&=NR_0,\qquad K=R_0(2B_0)^{J-1}.\label{Kbound}
\end{align}
For each $s=0,\ldots,J$, put
\begin{align}
 \Delta_s&=a_s^2,& C_s&=\Delta_s+(4\ell B_0+2)K,\label{Cs}\\
 M_s&=\Delta_s\sum_{i=0}^{m}
                   C_s^{i+1}\Delta_s^{m-i}A_i,&
 H_s(U)&=\det(I_n-UM_s).\label{Ms}
\end{align}
Finally define
\begin{equation}\label{Fformula}
 \boxed{\displaystyle
 F(T,U)=\prod_{s=0}^{J}
       \left(H_s(U)^2+\sum_{j=s+1}^{J}a_j^2\right).}
\end{equation}
All empty sums are zero. Every determinant, sum and product is finite,
and every entry belongs to $R$. This is an explicit coefficient
specification for an integer polynomial, even when expanding it would
be impractical.

\section{Why the element separates the permitted characters}
Fix a rational base specialization. Thus all matrix entries and coefficients
below are real rational numbers, although their eigenvalues may be complex.
Let $s$ be the largest index with $a_s\ne0$, taking $s=0$ if all
$a_1,\ldots,a_J$ vanish. Then
\[
 \chi(Z)=Z^{J-s}(Z^s+a_1Z^{s-1}+\cdots+a_s).
\]
Every eigenvalue of $\widetilde A_i$ has absolute value at most $B_0$:
each matrix entry has absolute value at most $R_0$, and the maximum row
sum bounds the spectral radius. Hence every root of $\chi$ has absolute
value at most $2B_0$.

If $\delta$ is a nonzero root of $\chi$, the product of its $s$ nonzero
roots, counted with multiplicity, has absolute value $|a_s|$. Therefore
\begin{equation}\label{lowerseparation}
 |\delta|\ge\frac{|a_s|}{(2B_0)^{s-1}}
       \ge\frac{a_s^2}{R_0(2B_0)^{J-1}}
       =\frac{\Delta_s}{K}=:\varepsilon.
\end{equation}
We used $|a_s|\le R_0$ and $2B_0>1$. When $s=0$ there are no nonzero
roots and the following separation statements are vacuous unless a
permitted character exists; such a character would itself give a
nonzero root through $\gamma_0$ and the extra zero eigenvalue. Indeed,
when $s=0$ all $\gamma_i$ have only zero character values; then
$H_0(U)=1$ and there is no permitted character.

Since $\Delta_s>0$, put
\[
             B=\frac{C_s}{\Delta_s}
               =1+\frac{4\ell B_0+2}{\varepsilon}>1.
\]
For any two different vectors in
\[
 \bigl\{(\sigma(\gamma_0),\ldots,\sigma(\gamma_m)):\sigma
                 \text{ a geometric character}\bigr\}\cup\{0\},
\]
let $h$ be their largest differing coordinate. Its difference has absolute
value at least $\varepsilon$, by \eqref{lowerseparation}; every coordinate
difference has absolute value at most $2B_0$. The leading term of their
weighted difference satisfies
\[
 \left|B^{h+1}\delta_h\right|\ge\varepsilon B^{h+1},
 \qquad
 \left|\sum_{i<h}B^{i+1}\delta_i\right|
        \le2\ell B_0B^h<\varepsilon B^{h+1}.
\]
Thus the weighted sums are distinct. The matrix in \eqref{Ms} is
multiplication by
\[
 \beta_s=\Delta_s^{m+2}\sum_{i=0}^{m}B^{i+1}\gamma_i.
\]
The scalar factor is nonzero and does not affect separation.

If $\sigma(g)=0$, every $\sigma(\gamma_i)$ is zero and
$\sigma(\beta_s)=0$. If $\sigma(g)\ne0$, the vector is nonzero and
$\sigma(\beta_s)\ne0$. Finally, two characters with nonzero guard have
different vectors unless they are the same character: recover their
generator values as $\sigma(\theta_i)=\sigma(\gamma_i)/\sigma(\gamma_0)$.
Consequently $\beta_s$ separates all permitted characters and sends
exactly the excluded characters to zero.

\section{Rationality and the coefficient strata}

\begin{proof}[Proof of Theorem~\ref{main}, equivalence]
For the actual last nonzero coefficient index $s$, suppose first that a
permitted character $\phi:\A_t\to\Q$ exists. Its value
$\phi(\beta_s)$ is a nonzero rational number. Choose
$u=1/\phi(\beta_s)$. Lemma~\ref{spectrum} gives $H_s(u)=0$;
all $a_j$ for $j>s$ vanish, so \eqref{Fformula} is zero.

Conversely, suppose $F(t,u)=0$ with $t,u$ rational. One factor of
\eqref{Fformula} is zero. It is a sum of rational squares, so its
$H_s(u)$ and all its tail coefficients vanish. If $a_s=0$, then
$\Delta_s=0$, $M_s=0$ and $H_s(u)=1$, a contradiction. Thus $s$ is
exactly the largest nonzero coefficient index. The determinant equality
also implies $u\ne0$ and supplies a geometric character $\sigma$ with
\[
                         \sigma(\beta_s)=1/u\in\Q^\times.
\]
This character has nonzero guard. Every Galois conjugate of $\sigma$
has the same $\beta_s$ value, because that value is rational. Separation
therefore makes every such conjugate the same character. In particular,
every generator value $\sigma(\theta_i)$ is fixed by the absolute
Galois group of $\Q$, hence belongs to $\Q$. The character takes values
in $\Q$ and has nonzero guard, as required.
\end{proof}

\begin{proof}[Proof of the degree bound]
The entries of $D_i$ have degree at most $d$, so
$\deg a_j\le jd\le Jd$. It follows that
\[
 \deg R_0\le2Jd,\quad
 \deg\Delta_s\le2Jd,\quad
 \deg K\le2J^2d,\quad
 \deg C_s\le2J(J+1)d.
\]
In the sum defining $M_s$, the largest bound comes from $i=m=\ell-1$.
Every entry of that matrix has degree at most
\[
 d_\beta=d+2Jd+2\ell J(J+1)d
          =d\{1+2J(1+\ell(J+1))\}.
\]
Hence $\deg H_s\le n(1+d_\beta)$. This bound also dominates the
degrees needed for the squared tail coefficients in \eqref{Fformula}.
There are $J+1$ factors, proving \eqref{degree}.
\end{proof}

\section{Clearing denominators without losing selected characters}
A useful input model is a triangular finite-root specification. At each
stage a root $\theta$ satisfies a monic polynomial with coefficients
rational in the ordinary base coordinates and preceding roots. Collect
its denominators in a common expression $D$, and include $D$ in the
nonzero guard at the selected character. If
\[
 \theta^e+\sum_{j=0}^{e-1}\frac{N_j}{D}\theta^j=0,
\]
replace the finite root by $\xi=D\theta$. The equation becomes the monic
relation
\[
                 \xi^e+\sum_{j=0}^{e-1}N_jD^{e-1-j}\xi^j=0.
\]
On the guarded locus this is reversible by $\theta=\xi/D$. Continue
in the triangular order, clearing the new denominators after each
substitution. A finite tower of monic polynomial extensions is free, with
basis the monomials whose exponent in each generator is below its stage
degree. The rank is the product of the stage degrees. All source guards
become a single polynomial guard after multiplying their numerators and
the recorded denominators. Thus Theorem~\ref{main} applies.

This procedure does not require a denominator to be nonzero at every
geometric character. The product guard selects the allowed characters.
Characters at which a cleared denominator vanishes remain in the finite
algebra and are rejected by the same separator. No extra ordinary
rational coordinate is introduced when a finite root is rescaled.

\section{Merging independent encoders}

\begin{corollary}\label{merge}
Suppose each of $k$ finite-root predicates has $b_i$ ordinary rational
base coordinates, a finite free algebra of rank $n_i$, and a polynomial
nonzero guard. Their conjunction is expressible with all ordinary base
coordinates and just one additional rational coordinate. Its finite
algebra can be taken to be the tensor product, of rank $\prod_i n_i$,
and its guard to be the product of the guards. No simultaneous
realization assertion beyond the original separate predicates is needed.
\end{corollary}
\begin{proof}
A character of the tensor product into $\Q$ is exactly a collection of
characters of its factors into $\Q$, agreeing on the rational base.
Its product guard is nonzero exactly when all original guards are nonzero.
Apply Theorem~\ref{main}. In particular, a tensor product such as
$\Q[\epsilon,\delta]/(\epsilon^2,\delta^2)$ presents no difficulty;
monogenicity of the full nonreduced algebra is never used.
\end{proof}

Applied to a source integer predicate with $m$ total rational coordinates,
comprising $m-1$ ordinary base coordinates and one final finite-root encoder, $k$ independent input
tests use $k(m-1)+1$ coordinates including the $k$ inputs. The mixed source
with eight integer unknowns and one nonnegative unknown then uses
\[
                    9(m-1)+1-1+3=9m-6.
\]
Here the nonnegative input is replaced by $s^2+t^2+u^2+u$, whose integral
values are exactly the nonnegative integers. For the ordinary-$abc$
interface first feed each finite-root source $X_i^2+1$. This expression is
positive and is integral exactly when $X_i$ is integral, so the external
zero-input factor is unnecessary. Thus the project's
25-, 27- and ordinary-abc 11-coordinate finite-root interfaces give
respective unrestricted-degree transfer counts
\[
                              219,\qquad237,\qquad93.
\]
These applications retain all the arithmetic premises of the relevant
integer predicates. They need a different determinant specification from
nine separate polynomials. The old numerical degree bound must not be
reused without checking \eqref{degree} for the new matrices.

Further identification of ordinary base coordinates requires a separate
simultaneous-completeness proof. Once such a proof is supplied, the same
encoder applies to the resulting joint triangular root algebra. For a
shared root one includes its monic equation once, so consistency of its
selected value across all blocks is imposed automatically.

\section{The mixed-integer source and its rational wrapper}
We use Sun's theorem \cite[Theorem 1.1(i)]{Sun}: for a fixed computably
enumerable nonrecursive set $W$, an integer-coefficient polynomial
$Q(x,z_1,\ldots,z_8,n)$ represents membership $x\in W$ by eight integer
unknowns and one nonnegative integer unknown. This classical theorem is
imported; the simultaneous integer certificates below are the new transfer.
The instance $x$ occurs in coefficients after specialization and is not
an additional existential variable.

\begin{lemma}\label{wrapper}
Put $L_+(s,t,u)=s^2+t^2+u^2+u$. Its integral values at rational arguments
are precisely the nonnegative integers, and each such integer has integer
arguments representing it.
\end{lemma}
\begin{proof}
Completing the square gives $L_+\ge-1/4$, so an integral value is nonnegative.
For $n\ge0$, the positive integer $4n+1$ is not of the form
$4^a(8b+7)$. The three-square theorem writes it as a sum of three integer
squares. Reduction modulo four shows exactly one of those integers is
odd. Write the other two as $2s,2t$ and the odd one as $2u+1$.
Then $4n+1=4s^2+4t^2+(2u+1)^2$, proving $n=L_+(s,t,u)$.
\end{proof}

If $\Phi_k$ is a joint certificate for $\mathbb Z^k$, then
\[
 Q(x,z_1,\ldots,z_8,L_+)^2+
 \Phi_9(z_1,\ldots,z_8,L_+,\mathbf w)^2=0
\]
has rational solutions exactly for $x\in W$. Substituting $L_+$ removes
one old input and introduces three coordinates. If $D$ bounds the degree
of $\Phi_9$ and $D_Q$ bounds that of $Q$, the final degree is at most
$4\max(D,D_Q)$. The explicit matrix bound \eqref{degree} supplies $D$.

\paragraph{Attribution and premises.}
The underlying source constructions include the notes of Gabe Conant with
Astra~6.2 and the subsequent project refinements. We use the full auxiliary
charts, four-contact scalars and twelve-contact ordinary-$abc$ detector
from v78, with the v88 normalized Bhargava--Burungale product poles.
Their elliptic-group, trace or effective abscissa-sum height, contact-height,
compatible prime-pattern and recurrence premises are retained
\cite{v19,v78,v88,BB}. The prime-pattern premise expressly
allows every fixed finite list of rational contact coefficients with
arbitrary finite avoidance and fixed divisibility. It is this stated
premise, rather than an unstated uniformity assertion, which is used for
the unions below. The simultaneous proofs and joint encoder do not
independently establish those arithmetic foundations.

\section{A joint certificate for finitely many integer coordinates}
\label{sec:joint-few-variable}

The unscaled point coordinates in the 25- and 27-coordinate sources
permit substantially more sharing than the quartic charts.  We use
the four-contact revision
\nolinkurl{papers/few_variables25/four_contact_scalar_20261119_v78.tex},
with the class-seven scalar and Bhargava--Burungale product-pole
replacement of
\nolinkurl{papers/few_variables25/bb_product_degree_transfer_20261123_v88.tex}.
We retain the full auxiliary point charts and the non-pole arithmetic
premises of those sources.  The two-detector version has an ineffective
sufficient fixed odd exponent; the three-detector version is the
specified member with exponent $L=143$.  Neither statement below
removes the retained elliptic, modular, or prime-pattern premises.
The published Bhargava--Burungale theorem supplies the pole converse
through the v88 transfer; the former native-pole realization premise
is not required.  No additional $abc$ hypothesis is used.

\begin{theorem}[Joint defining polynomials]
\label{thm:joint-few-variable}
For every fixed positive integer $k$ the following assertions hold
under the corresponding retained arithmetic premises.
\begin{enumerate}
\item For every sufficiently large fixed odd exponent $L$ there is a
finite integer-polynomial specification in $9k+16$ total rational
coordinates, including $k$ distinguished inputs, whose input
projection is exactly $\mathbb Z^k$.
\item There is a fully specified finite integer-polynomial
construction, at exponent $L=143$, in $11k+16$ total rational
coordinates with input projection exactly $\mathbb Z^k$.
\end{enumerate}
The sufficient exponent in the first assertion is the one required
by the single-input soundness theorem; it is not enlarged as $k$
increases.  The polynomials have a computable finite degree.  No
quartic degree bound, or retention of the old single-input degree
bound, is asserted.
\end{theorem}

\subsection{The point roles that can be shared}

Use the source curve
\[
 K=\mathbb Q(\sqrt2),\qquad D=11-6\sqrt2,\qquad
 E:y^2=x^3-Dx,
\]
and the source's fixed generator of $2E(K)$,
\[
 P=\left(\frac{57-28\sqrt2}{8},\frac{-660+417\sqrt2}{32}\right).
\]
The nine source point roles
are
\[
                     A,S,T,Z,H,U_A,U_S,U_T,U_Z.
\]
Each role has two retained rational coordinates.  A root role uses
the square-abscissa chart; an auxiliary role uses the full chart.
The finite chart generators are algebraic roots which are later
eliminated, and are not retained rational coordinates.

For input $x_i\in\mathbb Z$, the prepared witnesses use
\[
 n_i=x_i^2+1,\qquad A_i=[n_i]P,\qquad U_{A_i}=[hn_i]P.
\]
The remaining seven points can be common:
\[
 \begin{split}
 S&=[a]P,\quad T=[t]P,\quad Z=[b]P,\quad H=[h]P,\\
 U_S&=[ha]P,\quad U_T=[ht]P,\quad U_Z=[hb]P.
 \end{split}
\]
Thus only $A_i,U_{A_i}$ vary as point roles with $i$.  This includes
$x_i=0$, because its first root label is $n_i=1$, not zero.

The equations do not contain $a,t,b,h$ as extra quantified rational
coordinates.  They are integer labels used to construct the actual
elliptic points in the completeness proof.  Sharing seven points
means identifying their existing fourteen rational base coordinates
and their finite chart generators.

\subsection{Common scalar preparation}

For each input retain its own scalar
\[
 \begin{split}
 \beta_i&=1632x_i^2+2477,\qquad
 A_i^{\rm sc}=\beta_i(\beta_i-1),\qquad
 B_i^{\rm sc}=(\beta_i-2401)(32\beta_i+343),\\
 T_L^{\rm sc}&=\bigl(1224(67^{(L+1)/2}+1)\bigr)^2,\\
 \rho_i&=102^2 67^L\,
       \frac{7T_L^{\rm sc}A_i^{\rm sc}B_i^{\rm sc}}
            {(B_i^{\rm sc}+7T_L^{\rm sc}A_i^{\rm sc})^2}.
 \end{split}
\]
The v88 source proves $0<\rho_i<1/167$, square reduced denominator,
$\rho_i\in7\mathbb Q_2^{\times2}$ with $v_2(\rho_i)\ge10$,
the required finite preparation depths at $3,17$, and the fixed odd
valuation at 67.  No two of these scalars are identified.  The
17-adic square class need not be trivial: that condition belonged
to the replaced native-pole converse and is not used here.

\begin{lemma}[Simultaneous scalar preparation]
\label{lem:joint-few-scalar}
For a finite tuple of integer inputs, the source scalar preparation
can be performed with a common initial reference point $H$ and
a common integer $c\in6\mathbb Z_{>0}$, with $a=c+1$ and
$S=[a]P$.  The subsequent Lucas labels $t,b$ and the seven point
roles displayed above can consequently be common.
\end{lemma}
\begin{proof}
In the source, the six rational scalar components are the two
components of each of
\[
                   f(S-P),\qquad f(H),\qquad
                   z(S-P)/z(H),
\]
where $f(Q)=1/x(Q)$ and $z(Q)=-x(Q)/y(Q)$.  They depend on
$S-P=[c]P$ and $H$, and not on the input.  The input enters their
detector factors through the fixed number $\rho_i$.

Take the union of the finite preparation sets for all inputs.  It
includes the supports of the $\rho_i$, the fixed exceptional set,
and $2,3,17$.  Choose $H=[h_0]P$, with $64\mid h_0$, deep in the
formal groups at these finitely many places and in a small real
neighborhood of the identity, with nonzero local parameters.
For each of the finitely many inputs the source requires only
specified finite smallness tolerances for the six scalar components.
Use their minimum at each place.  The source construction first
fixes $H$ and then takes $[c]P$ still closer to the identity, with
$6\mid c$, so all these tolerances can be met at once.  At fixed
finite places this follows from the sufficiently deep formal
subgroup and continuity argument retained in the source.  At the
real places it follows from recurrence in the compact closure of
the cyclic orbit, with the finite component index included in the
progression.  Arbitrarily large positive such $c$ are available.

More explicitly, for every detector contact function $R$ the source
uses
\[
       G_\nu(R)=R^{\kappa_\nu}
                       \frac{R S_\nu}{1+R S_\nu},
\]
where $S_\nu$ is one of those six scalar components.  With each
$\rho_i$ fixed, the scalar-factor second summands tend to zero
as the $S_\nu$ tend to zero, uniformly for real $1/2<R<2$ and
for $R$ sufficiently close to 1 at the finitely many preparation
places.  The finite minimum of the required neighborhoods is
therefore still an admissible neighborhood.  In particular the
linear normalization denominators can all be kept nonzero.

Now set $a=c+1$ and take the source Lucas sequence
\[
 U_0=0,\quad U_1=1,\quad U_{j+1}=aU_j-U_{j-1},\qquad
 \ell=2(a-2),\quad b=U_\ell/(a-2),\quad t=U_{\ell+1}.
\]
The source congruences make $b,t$ positive odd integers and give
\[
                  t^2-a(a-2)tb+(a-2)^2b^2-1=0.
\]
They do not depend on the input.  All root points, scalars, and
comparison scales are fixed at this stage.  After those choices
the scalar preparation conditions are open in $H$.  Intersect the
finitely many allowed reference neighborhoods and preserve the
initial congruence class.  The intersection contains the initial
prepared point and is nonempty.
\end{proof}

There are finitely many further point and difference chart guards.
For fixed nonzero labels each multiplication or translation map
has only finitely many poles or forbidden point preimages.  Their
finite union can be avoided within the allowed reference
neighborhood.  Choose an actual allowed orbit point of index greater
than two outside that finite union, and shrink around it.  This
applies to all the $n_i$, including repeated values.  It supplies
one common neighborhood on which every raw
comparison function of every block is bounded.  In the full
auxiliary charts one may retain the original prepared subfamily in
$2E(K)$, exactly as in the single-input completeness proof.

\subsection{One contact parameter for every block}

Let $e=2$ for the general source and $e=3$ for the specified source.
Write
\[
 \Theta=\{0,-1/2,-1,1/2399,-32/407\},\qquad
 \Theta_j=(4j+2+\Theta)\cup(4j+4+\Theta),\quad 0\le j<e.
\]
Every element of every $\Theta_j$ lies strictly between 0 and 13.
The source arguments have the form
\[
                         t_{ij}=\alpha_{ij}\lambda_i,
\]
where $\alpha_{i0}=1$. Define
\[
 \Sigma_i=x(A_i)+x(S)+x(T)+x(Z)=U_i+\sqrt2 S_i,\qquad
 \mathcal T_i=2U_i,\qquad V_i=U_i+S_i.
\]
Here $S_i$ is a rational component, distinct from the common point $S$.
For $e=2$ the second scale is
$1+\rho_i\mathcal T_i^2$, and for $e=3$ the other two scales are
$1+\rho_i U_i^2$ and $1+\rho_i V_i^2$.  Here the trace quantities
are rational functions of the already fixed four root points.
Every scale is at least 1 on the prepared witnesses. Put
\[
 \delta_{ij}=t_{ij}+j+1,\qquad
 R_{ij}=\frac{\prod_{\theta\in\Theta_j}(t_{ij}-\theta)}{\delta_{ij}^{10}}.
\]
The original guards require $\delta_{ij}\ne0$.

In the new rational source make the fixed affine substitution
\begin{equation}\label{eq:joint-affine-contact}
                       \lambda_i=\lambda+100i,
                       \qquad 0\le i<k.
\end{equation}
Only $\lambda$ is retained.  Every old guard $\lambda_i\ne0$
is retained after substitution.

\begin{lemma}[The simultaneous contact family]
\label{lem:joint-few-contacts}
After the common scalar and root choices, the retained prime-pattern
and compatible annihilation supplier gives comparison witnesses
for all blocks with the substitution
\eqref{eq:joint-affine-contact} and a single final index $h$.
\end{lemma}
\begin{proof}
The contact positions as functions of $\lambda$ are
\[
                d_{ij\theta}=\theta/\alpha_{ij}-100i,
                            \qquad \theta\in\Theta_j.
\]
They are fixed rational numbers before prime selection.  For each
$i$ they lie in the interval $(-100i,-100i+13)$, so positions from
different blocks cannot coincide.  Within one block repeated
positions are merged with all incidences retained, as in the
single-input v78 construction.  Such incidences impose the same
successful elliptic comparison conditions.  Within each individual
contact function the ten positions are distinct.

Apply the retained prime-pattern theorem once to this finite
union, with
\[
                  \lambda=u/v,\qquad v=D_0^2w^2,
                  \qquad p_d=u-dv.
\]
Enlarge finite avoidance to include every scalar support, every
root-label numerator and denominator support, every scale support,
every nonzero contact difference, and every previously fixed
preparation modulus and guard denominator.  Include the common
contact denominator $D_0$ in that fixed finite set as well.  The
arbitrary fixed divisibility of $w$ makes $\lambda$ tend to infinity at all prepared
finite places.  Fixed translations by $100i$ do not change that
limit.  At these primes the fresh prime values force $u$ to be a
unit, so the prescribed negative valuations of $u/v$ hold.
The arbitrary real slope makes every $\lambda_i$ large
and positive.  Thus all source local and real limits hold
simultaneously.

The translated positions need not be positive.  We therefore check
the prime-order separation with both signs allowed.  Impose
\[
                 \lambda>4\max_d|d|,\qquad u\ge100.
\]
Every selected prime lies in $(3u/4,5u/4)$.  At either conjugate
of any selected prime the elliptic reduction group has even order,
and Hasse's bound gives
\[
 \#E(\mathbb F_p)\le p+1+2\sqrt p
       <\frac54u+1+\sqrt{5u}<\frac32u.
\]
The last inequality holds for $u\ge100$.  If a selected odd prime
divided any such even order, that order would be at least twice
the selected prime, hence strictly greater than $3u/2$, a
contradiction.  Every selected prime is therefore coprime to
every inherited reduction-order modulus and to the earlier
preparation modulus.

The retained compatible primary and conjugate annihilation
conditions give a common progression for $h$, preserving the
prepared reference congruence.  Add
\[
                            h\equiv0\pmod{\prod_d p_d}.
\]
This is compatible by the preceding coprimality.  The same
conditions work for every $n_i$, because all numerator points are
multiples of the same reference generator and their fixed nonzero
labels have been excluded from the selected primes.  There are
no input-dependent conflicting prescribed residues for $h$.
The source compact-component argument preserves the real return
on every compatible refinement.  Choose a final $h>2$ in this
progression and in the common allowed neighborhood.

At a selected prime of an individual detector its contact
function has valuation one, its $\rho_i$ and normalization
denominator are units, and the successful-slot point and comparison
depths are exactly those in the source proof.  All detector factors
are consequently one modulo the prime in that slot, and their
opposite inverse factors have residues zero or one.  Its reciprocal
trace is a unit.  The source scalar-normalization check at the fresh
primes also applies simultaneously: the rational components of
$f(H)$ are units, those of $f(S-P)$ are integral, and the components
of $z(S-P)/z(H)$ have valuation at most $-2$.  Hence no scalar
component has the forbidden valuation $-1$, and no linear scalar
normalization denominator vanishes there.

The factorization of each individual contact numerator is still
\[
 R_{ij}=\left(\frac{\alpha_{ij}}{v\delta_{ij}}\right)^{10}
         \prod_{\theta\in\Theta_j}
           \bigl(u-(\theta/\alpha_{ij}-100i)v\bigr).
\]
Its scale factor is a square and all its selected odd valuations
are positive.  The v88 parity proof, the finite preparations, and
the limits $R_{ij}\to1$, $W_{ij}\to1$ now apply to every actual
pole input
\[
                    e_{ij}=\frac{4\rho_iR_{ij}}
                              {\operatorname{Tr}(W_{ij}^{-1})^2}.
\]
Each has square reduced denominator.  At two it belongs to
$7\mathbb Q_2^{\times2}$ with depth at least ten, because $R_{ij}$
is close to one and the trace is close to two.  At the real place
$R_{ij}<2$ and the trace is greater than one, so
$0<e_{ij}<8\rho_i<1$.  These are precisely the prepared inputs
in the v88 Bhargava--Burungale transfer.  Extract the already
available square factor $e_{ij}=d_{ij}k_{ij}^2$ and use
\[
 Y_{ij}^2=(1-e_{ij})(v_{ij}^4-d_{ij}^2)
                         (s_{ij}^4-d_{ij}^2).
\]
The v88 pole lemma, proved from Bhargava--Burungale Lemma 2.2
and Theorem 2.3, supplies these witnesses independently for every
pair $(i,j)$, with Zariski-dense choices of $(v_{ij},s_{ij})$.
Thus the finitely many remaining nonidentically-zero pole guards
can be avoided.  Terminal zero roots are permitted.  No pole
parameters from different blocks have been identified.
\end{proof}

\subsection{The single finite-algebra encoding}

Retain one common finite chart algebra for each of the seven shared
point roles, two further point chart algebras for each input, and
$e$ terminal quadratic pole algebras for each input.  A point chart
has rank four.  The full finite algebra therefore has rank
\begin{equation}\label{eq:joint-few-rank}
                   n=4^{7+2k}2^{ek}
                     =2^{14+(4+e)k}.
\end{equation}
Use the source polynomial clearing of its rational coefficients
and retain the product of all guards.  Equivalently, any additional
matrix denominators can be cleared and included in the guard before
applying the guarded finite-algebra construction.  A guarded rational
character then gives a rational character of every individual source
block, with the seven common point characters agreeing.

Apply the general one-variable guarded rational-character encoder
to this finite algebra.  Its validity for arbitrary nilpotent
specializations is required: tensoring several terminal quadratic
algebras can produce more than one repeated-root direction.  It
would not suffice to assume that the tensor algebra is monogenic
or that all terminal roots are nonzero.  The general encoder adds
one rational coordinate and is equivalent to the existence of a
guarded rational character.  Its determinant and finite-product
recipe gives a finite integer-polynomial specification for each
fixed $k$ and fixed exponent $L$.

\begin{proof}[Proof of Theorem~\ref{thm:joint-few-variable}]
For soundness, a rational zero of the joint polynomial gives a
guarded rational character of the joint finite algebra.  Restricting
that character to each block, and undoing only the fixed affine
substitution for its slope, gives an original valid single-input
witness.  The original soundness theorem forces every input to
be an integer.  Its height estimates and scalar bounds are used
separately for each block, so their fixed sufficient exponent is
unchanged.

For completeness, Lemmas~\ref{lem:joint-few-scalar} and
\ref{lem:joint-few-contacts} give simultaneous rational source
witnesses for every integer tuple, including tuples with zero
or repeated entries.  Choose the same point-chart root values for
each common point.  These values and the independent pole-chart
roots define a guarded rational character of the joint algebra.
The general finite-algebra encoder supplies its one rational
encoding coordinate.

The complete count before this last coordinate is
\[
 \underbrace{k}_{\text{inputs}}+
 \underbrace{14}_{\text{common points}}+
 \underbrace{4k}_{\text{private points}}+
 \underbrace{1}_{\text{common slope}}+
 \underbrace{2ek}_{\text{pole parameters}}
                 =(5+2e)k+15.
\]
The single encoding coordinate gives $(5+2e)k+16$, proving the
two asserted counts.  For $k=1$ these specialize to the original
25 and 27 coordinates, as they should.
\end{proof}

\subsection{Transfer to fixed-variable undecidability}

\begin{corollary}[Ninety-nine and 117 variables]
\label{cor:99-117}
Under the retained premises of the two-detector general source,
solvability over $\mathbb Q$ of integer polynomial equations in
99 rational unknowns is undecidable.  Under the retained premises
of the specified three-detector source, a fully specified reduction
uses 117 rational unknowns.  Both assertions have a finite fixed
degree bound; this section makes no quartic claim.
\end{corollary}
\begin{proof}
Use the integer undecidability theorem with eight unrestricted
integer unknowns and one nonnegative integer unknown.  Set
\[
                    L_+(s,t,u)=s^2+t^2+u^2+u.
\]
Lemma~\ref{wrapper} shows that every nonnegative
integer has such a representation with integer $s,t,u$, and that
every integral value of $L_+$ at rational arguments is nonnegative.
In the joint certificate for $k=9$, substitute $L_+$ for its ninth
input coordinate.  This replaces one coordinate by three, and so
adds exactly two coordinates.  Combine its polynomial with the
specialized universal integer polynomial by the sum of their
squares.  Rational solvability is equivalent to the original
mixed-integer instance.  The two counts are
\[
                    9\cdot9+16+2=99,
             \qquad 11\cdot9+16+2=117.
\]
The universal polynomial has fixed degree.  The joint encoder also
has a computable finite degree for the fixed choice $k=9$; the
substitution and sum of squares therefore preserve a finite fixed
degree bound.  The input to the decision problem is encoded in
coefficients, not as another quantified unknown.
\end{proof}

\paragraph{Degree and effectiveness.}
For the general source the sufficient $L$ is ineffective.  The
undecidability conclusion is nevertheless an existence theorem for
a fixed polynomial reduction: choose one sufficiently large fixed
odd $L$.  The specified source uses $L=143$ and has no such
coefficient-parameter ambiguity.  The one-variable joint encoder
can have a much larger degree than the single-input encoder.  Its
rank at $k=9$ is $2^{68}$ for two detectors and $2^{77}$ for
three.  Reusing the old single-input degree or support bound after
this tensor construction would be unjustified.

\section{Sixty-nine variables under ordinary rational \texorpdfstring{$abc$}{abc}}
\label{sec:shared-abc}

The same joint finite-algebra method also improves the ordinary-$abc$
route.  Here the source is the twelve-contact eleven-coordinate
construction of v78, with its product-pole replacement in v88.
Its retained non-pole arithmetic premises, including ordinary
rational $abc$, are unchanged.  The published Bhargava--Burungale
theorem supplies the pole converse as in v88.  The sufficient fixed
odd shift $B$ remains ineffective.

\begin{theorem}[Joint ordinary-$abc$ certificate]
\label{thm:joint-abc}
For every positive integer $k$, under the retained ordinary-$abc$
source premises and for every sufficiently large fixed odd $B$,
there is a finite integer-polynomial specification in $7k+4$ total
rational coordinates, including $k$ inputs, whose input projection
is $\mathbb Z^k$.  Consequently polynomial solvability over
$\mathbb Q$ in 69 unknowns is undecidable under these premises.
The construction has a finite computable degree bound; it is not
a quartic construction.
\end{theorem}

\subsection{Removing the special zero-input branch}

The one-input source has three point roles $A,H,U$, one contact
parameter, two free pole parameters, and one finite-algebra encoder.
Its nonzero integer witnesses have
\[
                       A=[n]P,\qquad U=[n]H.
\]
The separate input factor in its final polynomial supplies the
integer zero.  Simply tensoring its nonzero-input algebras would
therefore not give a tuple certificate when some inputs vanish.

To avoid that problem, feed the $i$th source the rational expression
\[
                            y_i=X_i^2+1.
\]
This is positive at every rational base tuple, so the source's
zero-input factor can be omitted.  The equivalence
\[
                         X_i\in\mathbb Z
                    \quad\Longleftrightarrow\quad
                         X_i^2+1\in\mathbb Z
\]
follows by writing $X_i=r/s$ in lowest terms: if $s^2$ divides
$r^2+s^2$, then $s=1$.  Thus this substitution changes no intended
integer projection, uses no additional rational coordinate, and
allows every integer input tuple to use the nonzero source family.

Retain a private root point $A_i$ and a private numerator point
$U_i$ for each input, but identify the two rational base coordinates
and finite chart generators of the reference point $H$.  For
integer inputs choose
\[
                         A_i=[y_i]P,
                    \qquad U_i=[y_i]H.
\]
Every $y_i$ is a positive integer, including when $X_i=0$.

\subsection{A scaled common contact parameter}

The original source midpoint is
\[
 \beta_i=y_i^2+B,\qquad T_i=\operatorname{Tr}x(A_i),\qquad
                     a_i=\beta_i^4+T_i^4.
\]
It depends only on the input and the private root point $A_i$,
and is fixed before the reference and comparison-prime choices.
It is positive at every rational base tuple.  For $B\ge3$ one
has $a_i\ge B^4\ge81$.

In contrast with the previous general construction, the original
contact positions $\pm(a_i+c)$ have no uniform upper bound as
the input varies.  A fixed translation of the old contact parameter
alone would not justify disjointness.  Instead make the substitution
\begin{equation}\label{eq:abc-common-contact}
                         z_i=a_i(z+100i),
                           \qquad 0\le i<k.
\end{equation}
The multiplier $a_i$ is an already available expression, not an
extra coordinate.  The twelve source shifts are
\[
              \mathcal C=\{-9,-8,-7,-5,-4,-3,-2,-1,6,10,11,12\}.
\]
As functions of the common variable $z$, the contact positions are
\[
                       \pm(1+c/a_i)-100i,
                                  \qquad c\in\mathcal C.
\]
They lie in $(-100i-2,-100i+2)$ because $|c|\le12$ and
$a_i\ge81$.  Different input blocks therefore have disjoint
contact lists.  Within each block the 24 positions are distinct:
the shifts are distinct and the two signs cannot collide since
$a_i>12$.  Every position is fixed before the prime supplier is
invoked.

\subsection{Joint reference and pole preparation}

For an integer tuple first fix all $A_i=[y_i]P$ and hence all
$T_i,a_i$.  The v78 two-adic trace lemma gives
\[
 y_i^2T_i-57/4\in16\mathbb Z_2,\qquad
 v_2(T_i)=-2-2v_2(y_i),\qquad T_i\in\mathbb Q_2^{\times2}.
\]
Consequently $a_i=(y_i^2+B)^4+T_i^4$ is a two-adic square.
This is used only for prepared integer witnesses, not as an
assertion about all rational source tuples.

Take the union of the finite preparation sets for all input blocks,
including the supports of the $a_i$, root multipliers, and all
fixed denominators.  The reference neighborhoods and congruence
requirements are those in the original multiplication family
$U_i=[y_i]H$.  For the finitely many fixed multiplication and
translation maps, forbidden chart values and nonidentically-zero
raw-test zeros form a finite union.  Choose one allowed reference
orbit point outside this union and a common neighborhood around
it.  Raw tests which vanish identically cause no difficulty: their
positive detector factors are exactly one.  All other raw tests
can be bounded and kept away from zero in this common neighborhood.

Apply the retained prime-pattern theorem once to the union of the
normalized contact lists in \eqref{eq:abc-common-contact}, with
all finite denominator, avoidance, and approximation conditions
fixed beforehand.  The two-sided prime-cluster argument used in
Lemma~\ref{lem:joint-few-contacts} applies without change to these
rational positions.  It makes every selected prime coprime to
every selected reduction order.  The compatible index refinement
can therefore impose the required primary and opposite formal
conditions simultaneously on one $H=[h]P$.  Multiplication by
each fixed $y_i$, a unit at the selected primes, gives the source
successful comparison depths for that block.  The compact-component
recurrence preserves the common real neighborhood.

The fixed-place limits also survive the scaled substitution.
For the source derivative quotient $u_i$ and its normalization
$w_i=u_i/(1+u_i)^2$, one has
\[
 u_i\sim z_i^2/a_i=a_i(z+100i)^2,
          \qquad w_i\sim\frac{1}{a_i(z+100i)^2}
\]
at real and prepared finite infinity.  The fixed nonzero multipliers
$a_i$ only change the finitely many prescribed approximation depths.
At two their square classes are trivial by the preceding trace
lemma, so arbitrarily large negative $v_2(z)$ gives the same
positive even depth and square class as the source.  At seven
and every other prepared place the union of fixed supports and
sufficiently deep approximation gives the original nonnegative
odd-valuation conditions.

The positive balanced detectors retain $W_i\to1$ and reciprocal
trace $V_i\to2$ on these neighborhoods.  At every freshly
selected contact prime their successful-slot factors are one
modulo the prime and their opposite inverse factors have residues
zero or one, so the trace is a unit.  The unchanged source
all-prime argument gives square reduced denominators for the
actual inputs
\[
 e_i=\frac{28u_i}{(1+u_i)^2V_i^2}
    =7u_i\left(\frac{\operatorname{Nm}(W_i)}
                       {(1+u_i)(W_i)_0}\right)^2.
\]
The real preparation gives $0<e_i<1$, and the two-adic preparation
gives the square class seven and positive even depth at least four.
Thus the v88 Bhargava--Burungale product pole applies separately
with $d_i=7u_i$:
\[
             Y_i^2=(1-e_i)(v_i^4-d_i^2)(s_i^4-d_i^2).
\]
Each pair $(v_i,s_i)$ remains private.  Independent Zariski-dense
choices avoid the finitely many nonidentically-zero pole guards.
The terminal roots are allowed to vanish, as required by the v88
construction.

\subsection{Polynomial construction and count}

Keep one finite point algebra for the common $H$, two private point
algebras per input, and one terminal quadratic pole algebra per
input.  Their total rank is
\[
                         4^{1+2k}2^k=2^{2+5k}.
\]
Apply the general guarded finite-algebra encoder once to their
product guard, allowing all terminal nilpotent specializations.
Soundness follows by restricting a rational character to each
source block: it forces $y_i\in\mathbb Z$ and hence
$X_i\in\mathbb Z$.  The simultaneous construction above gives
completeness for every integer tuple, including zero and repeated
entries.  The old input factors are not multiplied into this joint
polynomial; their arguments are the everywhere-positive $y_i$.

The coordinate count is
\[
 \underbrace{k}_{\text{inputs}}+
 \underbrace{2}_{\text{reference}}+
 \underbrace{4k}_{\text{private points}}+
 \underbrace{1}_{\text{contact}}+
 \underbrace{2k}_{\text{pole parameters}}+
 \underbrace{1}_{\text{encoder}}=7k+4.
\]
For $k=9$, replacing the nonnegative integer input of the
mixed-integer universal equation by the three-coordinate wrapper
$s^2+t^2+u^2+u$ adds two coordinates, giving
\[
                              7\cdot9+4+2=69.
\]
Combining the universal equation and the joint certificate by a
sum of squares proves Theorem~\ref{thm:joint-abc}.  All substitutions,
finite multiplication tables, determinants, and products have
finite explicitly prescribed sizes, so the resulting degree is
fixed and computably bounded.  Its value is not the old single-input
degree bound.  At $k=9$ the joint algebra has rank $2^{47}$.

\section{Explicit rational source functions}
\label{sec:shared-source-functions}

For completeness, the rational functions used in the simultaneous
proofs are recorded here.  These formulas are the v78 sources with
the v88 scalar and product-pole changes.  All displayed divisions
carry their denominator guards; no denominator is discarded when
the equations are cleared.

The fixed point is
\[
 P=\left(\frac{57-28\sqrt2}{8},
         \frac{-660+417\sqrt2}{32}\right)
       \in E\bigl(\mathbb Q(\sqrt2)\bigr),
 \qquad E:y^2=x^3-(11-6\sqrt2)x.
\]
A square-abscissa root chart retains $h_0,h_1\in\mathbb Q$ and
uses
\[
 h=h_0+h_1\sqrt2,\qquad x=h^2,\qquad
 j^2=h^4-(11-6\sqrt2),\qquad y=hj.
\]
A full auxiliary chart retains $h_0,h_1$ and instead uses
\[
 x=h_0+h_1\sqrt2,\qquad
 j^2=x^3-(11-6\sqrt2)x,\qquad y=j.
\]
In either case, writing the right side of the square equation as
$A+B\sqrt2$, its finite rank-four algebra has basis $1,U,V,W$,
with $j=(U+V\sqrt2)/2$ and multiplication table
\[
 \begin{array}{lll}
 U^2=2A+2W,& V^2=A-W,& UV=2B,\\
 UW=AU-2BV,& VW=BU-AV,& W^2=A^2-2B^2.
 \end{array}
\]
These polynomial formulas define the free rank-four algebra at
every base specialization.  The source guards select its allowed
rational point characters.  In particular, the matrix construction
does not introduce a new rational variable for each of $U,V,W$.

\subsection{The general two- and three-detector sources}

For input block $i$, set
\[
 (Q_{i0},Q_{i1},Q_{i2},Q_{i3},Q_{i4})=(P,A_i,S,T,Z),
\]
\[
 (U_{i0},U_{i1},U_{i2},U_{i3},U_{i4})
                    =(H,U_{A_i},U_S,U_T,U_Z),
 \qquad V_{ir}=U_{ir}-Q_{ir}.
\]
The differences are evaluated by the guarded affine group formulas.
Write $\overline{\phantom{x}}$ for conjugation in
$\mathbb Q(\sqrt2)$, $f(Q)=1/x(Q)$, and $z(Q)=-x(Q)/y(Q)$.
Define
\[
 L_{ir}=\frac{z(U_{ir})}{z(H)},\qquad
 L'_{ir}=\overline{\frac{z(V_{ir})}{z(H-P)}}\quad(1\le r\le4).
\]
In order, the sixteen raw comparison functions are
\[
 \begin{gathered}
 (f(U_{ir}))_{r=0}^4,\qquad
 (\overline{f(V_{ir})})_{r=0}^4,\\
 L_{ir}-L'_{ir}\quad(1\le r\le4),\qquad
 L_{i1}-(X_i^2+1),\qquad
 \mathcal F(L_{i2},L_{i3},L_{i4}),
 \end{gathered}
\]
where
\[
          \mathcal F(u,v,w)=v^2-u(u-2)vw+(u-2)^2w^2-1.
\]
Call them $C_{ir}$, $1\le r\le16$.  Their required orders are
$a_r=3$ for $r\in\{12,13,14,16\}$ and $a_r=1$ otherwise.

The six common rational scalar components $S_\nu$ are, in order,
the coefficients in the basis $1,\sqrt2$ of
\[
                         f(H),\qquad f(S-P),\qquad
                         z(S-P)/z(H).
\]
Their orders are $\kappa_\nu=1$ for $1\le\nu\le4$ and
$\kappa_5=\kappa_6=3$.

Write
\[
 x(A_i)+x(S)+x(T)+x(Z)=\mathcal U_i+\sqrt2\mathcal S_i,
 \qquad\mathcal T_i=2\mathcal U_i,\qquad
 \mathcal V_i=\mathcal U_i+\mathcal S_i.
\]
The two-detector scales are
$\alpha_{i0}=1$, $\alpha_{i1}=1+\rho_i\mathcal T_i^2$.
The three-detector scales are
$\alpha_{i0}=1$, $\alpha_{i1}=1+\rho_i\mathcal U_i^2$,
$\alpha_{i2}=1+\rho_i\mathcal V_i^2$.
With the shifted common slope, put
\[
 t_{ij}=\alpha_{ij}(\lambda+100i),\qquad
 \delta_{ij}=t_{ij}+j+1,\qquad
 R_{ij}=\frac{\prod_{\theta\in\Theta_j}(t_{ij}-\theta)}
                  {\delta_{ij}^{10}},\qquad q_{ij}=\rho_iR_{ij},
\]
and
\[
       G_{ij\nu}=R_{ij}^{\kappa_\nu}
                   \frac{R_{ij}S_\nu}{1+R_{ij}S_\nu}.
\]
The positive field detector is exactly
\[
 \begin{split}
 W_{ij}={}&\prod_{r=1}^{10}
       \left[1+\left(\frac{C_{ir}}{\delta_{ij}^2q_{ij}}\right)^2\right]\\
 &{}\cdot\prod_{r=11}^{16}
       \left[1+\left(\frac{C_{ir}^2}
                         {\delta_{ij}^2q_{ij}^{2a_r-1}}\right)^2\right]
 \cdot\prod_{\nu=1}^{6}
       \left[1+\left(\frac{G_{ij\nu}^2}
                              {q_{ij}^{2\kappa_\nu-1}}\right)^2\right].
 \end{split}
\]
The reciprocal trace is $\operatorname{Tr}(W_{ij}^{-1})$ and the
actual pole input is $e_{ij}=4q_{ij}/\operatorname{Tr}(W_{ij}^{-1})^2$.
The complete v88 circuit specifies its already available square
extraction $e_{ij}=d_{ij}k_{ij}^2$ and the guarded normalization of
$d_{ij}$.  Changing that representative by a guarded square changes
only the two private pole parameters and its finite terminal root,
as proved in v88; it adds no rational coordinate.

\subsection{The ordinary-$abc$ source}

Here the private root and numerator points are $A_i,U_i$, with
common $H$.  Set
\[
 r_i=z(U_i)/z(H),\qquad
 r_i'=\overline{z(U_i-A_i)/z(H-P)}.
\]
The six raw tests are
\[
 f(H),\quad f(U_i),\quad \overline{f(H-P)},\quad
 \overline{f(U_i-A_i)},\quad r_i-r_i',\quad r_i-y_i,
                 \qquad y_i=X_i^2+1.
\]
Let $C_{i1},\ldots,C_{i4}$ be the first four tests themselves and
let $C_{i5},C_{i6}$ be the squares of the last two tests.
For the twelve shifts $\mathcal C$ stated in the proof, define
\[
 G_0(t)=\prod_{c\in\mathcal C}(t+c),\qquad
 F_i=G_0(a_i-z_i)G_0(a_i+z_i),\qquad
 H_{a,i}=\frac12\frac{\partial F_i}{\partial a_i},
\]
\[
 u_i=-\frac{12F_i}{H_{a,i}},\qquad
 z_i=a_i(z+100i),\qquad
 R_{ir}=\frac{C_{ir}^2}{u_i(1+u_i^2C_{ir}^2)},\qquad
 W_i=\prod_{r=1}^6(1+R_{ir}^2).
\]
The partial derivative is taken with $z_i$ formally held fixed,
before making the common-parameter substitution.  This order is
essential: differentiating the substituted expression as a function
of $a_i$ would define a different detector.
Finally
\[
 V_i=\operatorname{Tr}(W_i^{-1}),\qquad
 e_i=\frac{28u_i}{(1+u_i)^2V_i^2},\qquad d_i=7u_i.
\]
These are the v88 source functions whose terminal product-pole
equation is used in the joint construction.  Zero raw tests yield
the factor one, and zero terminal roots remain allowed.


\section{Explicit degree bounds and reproducible specifications}
\label{sec:joint-degree}
The matrix bound in Theorem~\ref{main} can be evaluated without constructing
matrices of rank $2^{68}$ or larger. Here is a reproducible bound on every
entry, including the effects of all denominator clearings.

Give every retained rational coordinate weight one and every fixed
coefficient parameter weight zero. A coefficient circuit is evaluated in
multiplication operators, with addition taking the maximum of degree bounds,
multiplication taking their sum, and a nonnegative integral power multiplying
the bound by its exponent. For a point block whose polynomial coefficients
$A,B$ have degrees at most $a$, the multiplication table in
Section~\ref{sec:shared-source-functions} shows that multiplication by
$U,V,W$ has entry degrees at most $a,a,2a$, respectively. All its $A,B$
depend only on retained coordinates; this dependency is checked in every
stored block. The table is independently verified by symbolic $4\times4$
matrix identities.

For a monic terminal stage $q^2+c=0$, the matrix of multiplication by $q$
on the two-term basis has entries $0,1,-c$. Thus its degree is bounded by
the already evaluated operator-degree bound for $c$. No terminal defining
coefficient depends on a terminal root. Extending a preceding multiplication
matrix to this stage makes block-diagonal copies, so its bound is unchanged.
Tensor products of the independent point blocks have the same property.
These observations prove the circuit degree induction, even at nilpotent
specializations. They do not rely on numerical sampling.

Let $g$ be the product of all retained guard numerators. Summing their
operator-degree bounds gives $d_g$. For the separating list we need only
$U,V$ from each point block, since $W=U^2/2-A$ on every rational fiber;
all terminal roots are retained in that list. The elements $U,V$ generate
the point algebra after extending constants from $\mathbb Z$ to $\mathbb Q$,
as allowed in Theorem~\ref{main}. If $d_\theta$ is the largest degree bound
for these generators, every $g\theta$ has degree at most $d=d_g+d_\theta$.
The exact ledger is
\begin{center}
\begin{tabular}{lrrrrr}
\toprule
Route & Rank $n$ & $\ell$ & $d_g$ & $d_\theta$ & $d$\\
\midrule
General & $2^{68}$ &69&2201890&100232&2302122\\
Specified & $2^{77}$ &78&4026046&100232&4126278\\
Ordinary $abc$ & $2^{47}$ &48&3103092&206888&3309980\\
\bottomrule
\end{tabular}
\end{center}
Substitution in \eqref{degree} gives the joint-certificate degree $D$.
The mixed-domain source can be chosen from the explicit construction in
\cite[Theorem~3, Definition~5.3 and Section~6.1]{BDHMS}. After substituting
$L_+$, its degree is at most
\[
\eta(\nu,\delta)=15616+233856\delta+
 (233952\delta+467712\delta^2)(2\delta+1)^{\nu+1}.
\]
We use its proved universal pair $(\nu,\delta)=(58,4)$, not the smaller
pair whose full classical proof is not supplied. Squaring the source and
the substituted joint certificate gives the safe bound
\[
                  \deg G\le\max\{2\eta(58,4),4D\}.
\]
The resulting numerical bounds are
\begin{center}
\begin{tabular}{lrr}
\toprule
Route & Rational unknowns & Degree strictly less than\\
\midrule
General &99&$2^{526}$\\
Specified &117&$2^{591}$\\
Ordinary $abc$ &69&$2^{377}$\\
\bottomrule
\end{tabular}
\end{center}
These are explicit upper bounds for fixed-degree families, not claims
that the displayed degrees are attained or optimal. The old independent
227-, 245- and 101-variable transfers have a much smaller common degree
bound $2\eta(58,4)$ and remain useful variable--degree tradeoffs.

The complete directed acyclic coefficient circuits and coordinate maps are
in \nolinkurl{undecidability/joint_sources/}, with one compressed JSON
specification and a certificate for each route. Together with the displayed
finite multiplication tables and the determinant formula \eqref{Fformula},
these specify each polynomial finitely and exactly; the enormous final
polynomials have not been expanded. The verifier
\nolinkurl{undecidability/verify_joint_sources.py} rebuilds all three
specifications byte for byte, checks every dependency needed above, and
recomputes all degree bounds with exact integers. Four separate exact
examples check the encoder on a quadratic field, a nonmonogenic algebra
with two nilpotent generators, a nilpotent guard excluding every character,
and a denominator which vanishes only at excluded characters.

\paragraph{Defining members and algebraic fixtures.}
The stored general-source value $L=143$ is an algebraic fixture, not a
certified sufficient specialization for the two-detector theorem. For that
route one must first select any sufficiently large fixed odd $L$ and run
the original v88 builder at that value, then apply the same sharing map.
The stored ordinary-$abc$ coefficient $B=3$ is likewise only a fixture;
the theorem requires a sufficiently large fixed odd $B$. The specified
117-variable route uses the certified source choice $L=143$.

The degree bounds hold uniformly over these fixed coefficient parameters:
$L$ occurs in coefficient constants such as $67^L$ and $T_L$, not in the
exponents of rational coordinates; $B$ is a coefficient parameter in the
same fixed circuit. Changing either parameter therefore leaves the
operator-degree induction valid. A smaller fixture does not certify the
arithmetic projection, and finite verification is not asserted to do so.
For the ineffective routes the conclusion is the existence of a fixed
finite defining coefficient list. Once such a list is fixed, its use in
the reduction is computable; an explicit numeric implementation must use
an actual defining member rather than the fixture.

\begin{thebibliography}{9}
\bibitem{BDHMS} J. Bayer, M. David, M. Ha{\ss}ler, Y. Matiyasevich and
D. Schleicher, \emph{Diophantine equations over the integers: universal
bounds and parallel formalization}, version~3, 22 October 2026.
\url{https://arxiv.org/abs/2506.20909v3}.
\bibitem{Sun} Z.-W. Sun, \emph{Further results on Hilbert's Tenth Problem},
Science China Mathematics \textbf{64} (2021), 281--306.
\url{https://arxiv.org/abs/1704.03504}, Theorem~1.1(i).
\bibitem{v19} \(\mathbb Z\mathbb Q\) Diophantine Project,
\emph{A defining polynomial in twenty-five variables}, v19 (2026).
\nolinkurl{papers/few_variables25/few_variables25_merged_20261110_v19.tex}.
\bibitem{v78} \(\mathbb Z\mathbb Q\) Diophantine Project,
\emph{Four-contact scalar reduction}, v78 (2026), and the twelve-contact
ordinary-$abc$ construction of the same revision.
\nolinkurl{papers/few_variables25/four_contact_scalar_20261119_v78.tex};
\nolinkurl{papers/conditional_few_variables/twelve_contact_first_scale_20261119_v78.tex}.
\bibitem{v88} \(\mathbb Z\mathbb Q\) Diophantine Project,
\emph{Product poles and lower degrees in the few-variable constructions},
v88 (2026).
\nolinkurl{papers/few_variables25/bb_product_degree_transfer_20261123_v88.tex}.
The precise retained non-pole premises and the normalized BB pole
completion are stated there.
\bibitem{BB} M. Bhargava and A. Burungale, the square-denominator and
product-pole results audited in the v88 source,
\url{https://arxiv.org/abs/2611.16130v2}, Lemma~2.2 and Theorem~2.3.
\end{thebibliography}

\end{document}
